\section{From bounded patterns to spherical arrangements}
\label{sec:planar}

Proposition~\ref{prop:bounded-pattern} controls systems of bounded size.
We now show that any spherical arrangement of the kind needed for the
topological argument contains such a system.  The number of disks in the
arrangement is unrestricted.  Throughout this section, probability is
taken in the distribution conditioned on girth at least $L$.

\begin{definition}[Reduced spherical arrangement]
\label{def:arrangement}
A reduced spherical arrangement over $\Gamma$ consists of the following
data on an oriented sphere.
\begin{enumerate}
\item There are finitely many pairwise disjoint piecewise linear closed
disks.  Each disk boundary carries a positive-length cyclically immersed
closed path in $\Gamma$, with the possible exception of one boundary.
That exceptional boundary carries a positive-length immersed interval
path starting at $x_A$ or $x_B$: its two endpoints are represented at one
marked boundary point, called the \emph{break}.  No immersion condition
is imposed across that break.  At least one ordinary cyclically immersed
boundary is present.
\item The paths are read in the boundary orientations induced by the
oriented disks.  Their edge occurrences are paired in distinct twos by
piecewise linear arcs whose interiors lie outside the disks.  The arcs
are mutually disjoint, and each occurrence contains exactly one arc end
in its interior.
\item Paired occurrences have inverse signed letters in these boundary
orientations, and their underlying undirected edges of $\Gamma$ are
distinct, as required by Equation~\ref{eq:edge-reduction}.
\end{enumerate}
Occurrences remain distinct even when the paths revisit the same graph
edge.  The arc graph may have loops, parallel edges, or several connected
components.
\end{definition}

\subsection{Short closures}

Fix a deterministic constant $d$ such that the event
\[
 \operatorname{diam}(\Lambda)\le dL
 \quad\text{for every connected component $\Lambda$ of $\Gamma$}
 \tag*{\textnormal{(diameter event)}}
\]
has probability $1-o(1)$, as supplied by Lemma~\ref{lem:diameter}.
We first work on this event.

\begin{lemma}[Closing a segment without cancellation]
\label{lem:short-closure}
There is a fixed integer $D$, independent of $n$ and of the segment,
such that every positive-length immersed segment in $\Gamma$ extends to
a cyclically immersed closed path by appending at most $DL$ edges.
The original segment is retained without cancellation.
\end{lemma}

\begin{proof}
Appending a geodesic return may cancel the ends of the given segment.
We prevent this by inserting short based loops with prescribed first
and last edge exclusions.  To construct such a loop at a vertex $z$,
delete at most two specified incident undirected edges from its
side of $\Gamma$.  For sufficiently large $n$, girth at least $L\ge3$
excludes graph loops and parallel edges; after the deletion every vertex
still has degree at least $(q+1)-2\ge3$.  Put
\[
 k=\lceil\log_2 n\rceil+1.
\]
From $z$ in the remaining graph there are at least $3\cdot2^{k-1}>n$
non-backtracking walks of length $k$.  That graph has at most $n$
vertices, so two distinct such walks $P,Q$ end at the same vertex.
Reduce $PQ^{-1}$ as a linear edge path.  This removes just their common
terminal segment.  Since $P$ and $Q$ have equal length and are distinct,
neither walk is entirely removed.  The result is a nonempty reduced
based loop of length at most $2k$, whose first and last edges avoid both
deleted edges.  The loop need not be cyclically immersed at $z$; only
its linear reduction will be used.

Let $S$ be the given segment, running from $a$ to $b$, and let $R$ be a
geodesic from $b$ to $a$, of length at most $dL$.  If $R$ is nonempty,
choose a based loop $J_b$ at $b$ by excluding the last edge of $S$ and
the first edge of $R$.  Choose $J_a$ at $a$ by excluding the last edge
of $R$ and the first edge of $S$.  Then
\[
 S J_b R J_a
\]
has no backtracking at any join, including its cyclic last/first join.
If $R$ is empty, so that $a=b$, use just one based loop $J_a$, excluding
the last and first edges of $S$; the path $SJ_a$ is cyclically immersed.
The immersion $\Gamma\to F$ also makes every one of these joins reduced
in signed letters.

For large $n$ we have $L\ge(c/2)\log n$ and $L\ge1$, whence
\[
 k\le \left(\frac{2}{c\log2}+2\right)L.
\]
Thus any integer
\[
 D\ge d+\frac{8}{c\log2}+8
\]
bounds the added length $|R|+4k$ in units of $L$, and covers the
empty-return case as well.
\end{proof}

\subsection{Turning planar pairings into intervals}

The next bound concerns only the original paired occurrences; edges
added to close segments will always be left unpaired.

\begin{lemma}[A linear bound on comparison intervals]
\label{lem:planar-intervals}
Suppose an arrangement with $N$ disk boundaries is divided into $N'\ge N$
pieces by cutting each original path into consecutive segments, without
crossing an exceptional break.  A path left intact counts as one piece.
Then its occurrence pairing can be specified by at most $6N'$ pairs of
intervals.  Both intervals in every pair lie wholly in their designated
pieces, are distinct as occurrence intervals, and compare inverse
strings in reversed boundary order.
\end{lemma}

\begin{proof}
Thicken the disks and pairing arcs to a regular disk-and-band
neighborhood $R$ of their planar multigraph.  Let $E$ denote its number
of arcs.  A \emph{vertex gap} is the boundary arc between consecutive
pairing-arc ends in the cyclic order around one disk.
For each connected complementary compact surface $Q$, let
$l_Q$ be the number of band sides in its boundary.  Equivalently, $l_Q$
counts the intervening vertex gaps, with incidences counted separately.
Every vertex has positive degree, so every boundary component of $R$
meets a band side, and $l_Q\ge1$.

Call a vertex gap \emph{good} if it belongs to a complementary disk
digon whose two sides come from two distinct arcs.  The two gaps across
such a digon are both good. If consecutive occurrences at one corner
are $e,f$ and their partners are $e',f'$, the opposite corner has
coherent order $f',e'$. Thus
$(e,f)\leftrightarrow(f',e')$ is an involution on good gaps. Here the
symbols denote occurrences, not underlying graph edges.
Good gaps therefore continue an inverse-word comparison.  We first
bound the remaining gaps, where this continuation can fail, before
making the required indexing and splitting cuts.

The neighborhood retracts to the multigraph, so $\chi(R)=N-E$.
Gluing the complementary surfaces to it along boundary circles gives
\[
 \sum_Q\chi(Q)=2-N+E,
 \qquad \sum_Q l_Q=2E.
\]
Consequently
\begin{equation}
 \sum_Q\bigl(l_Q-2\chi(Q)\bigr)=2N-4.
 \label{eq:planar-euler}
\end{equation}
This calculation uses neither connectedness of the multigraph nor
simple connectivity of its complementary regions.

We bound the non-good gaps, including the possible degenerate faces,
as follows.
\begin{itemize}
\item A disk monogon has one gap between occurrences paired by its
single arc.  Those consecutive letters are inverse, so immersion
forbids the gap unless it is the exceptional break.  Thus there are
at most $M\le1$ monogons, contributing $-M$ to
Equation~\ref{eq:planar-euler}.
\item A disk digon traversing the same arc twice must turn immediately
back along that arc at an endpoint.  No other arc end can intervene in
the cyclic order there, so that endpoint has degree one.  Different
such faces use different degree-one gaps.  There are at most $N$ such
faces, and hence at most $2N$ gaps of this type.  This deliberately
allows a coarse bound, including isolated single-edge components.
\item Every connected non-disk complementary surface is planar with
Euler characteristic at most zero.  Its length is therefore at most
its contribution $l_Q-2\chi(Q)$.  For a disk region of length $l_Q\ge3$,
we instead have $l_Q\le3(l_Q-2)$.
\end{itemize}
All contributions except those of monogons and digons are positive.
Their total is $2N-4+M$, by Equation~\ref{eq:planar-euler}.  If $B_{\mathrm{gap}}$ is
the number of non-good gaps, these observations give
\begin{equation}
 B_{\mathrm{gap}}\le M+2N+3(2N-4+M)\le8N.
 \label{eq:bad-gaps}
\end{equation}

Cut at every non-good gap.  Mark also one indexing gap for each original
path, using the exceptional break when present, and all gaps used in
the prescribed splitting.  There are at most $N+N'$ additional marked
gaps: a split cyclic boundary into $m$ pieces needs at most $m$ splitting
gaps, and a split interval needs no more.  Whenever a marked gap is
good, cut both it and its opposite good gap.  The resulting cut set is
stable under the good-gap involution and has cardinality $J$ satisfying
\[
 J\le B_{\mathrm{gap}}+2(N+N')\le12N'.
\]

Every original boundary now has at least one cut.  Its remaining
occurrence-adjacency components are therefore linear runs, all lying
inside their prescribed pieces.  The occurrence pairing reverses
every uncut adjacency.  Stability of the cut set implies that it maps
each whole maximal run onto another whole maximal run in reverse
order.  A run cannot map to itself: an odd-length run would have a fixed
occurrence, while an even-length run would pair its two central
consecutive occurrences and give inverse adjacent letters.  The latter
contradicts internal immersion; the exceptional break was cut and
cannot occur at the center of such a run.

There are exactly $J$ runs, since each original boundary is a circle
with a nonempty cut set.  They pair in distinct twos, giving at most
$J/2\le6N'$ interval pairs.  Inverse letters and reversed order are
exactly the signed comparison convention of
Proposition~\ref{prop:bounded-pattern}.
\end{proof}

\subsection{Separating a planar graph into bounded components}

\begin{lemma}[Recursive planar separation]
\label{lem:planar-separation}
Put
\[
 C_{\mathrm{sep}}=\frac{2\sqrt2}{1-\sqrt{2/3}}.
\]
For every integer $K\ge1$, a planar multigraph with $m$ vertices has a
set of at most $C_{\mathrm{sep}}m/\sqrt K$ vertices whose deletion leaves components
of size at most $K$.
\end{lemma}

\begin{proof}
Remove loops and merge parallel copies of edges.  These operations
preserve connected components, including after any vertex deletion.
The planar separator theorem of Lipton and Tarjan
\cite[Corollary~2, p.~183]{LiptonTarjan1979} supplies, in any simple planar graph with $s$
vertices, at most $2\sqrt{2s}$ vertices whose deletion leaves components
with at most $2s/3$ vertices.  Apply this theorem recursively to each
component with more than $K$ vertices.

For a recursive call on $s>K$ vertices, charge $2\sqrt2/\sqrt s$ to
each of those vertices.  The total charge at this call is
$2\sqrt{2s}$ and bounds the number deleted.  Fix one original vertex
and list the successive component sizes containing it before it is
deleted or recursion stops.  Consecutive sizes decrease by a factor
at most $2/3$.  Reading the list backwards from its last size greater
than $K$, its total charge is at most
\[
 \frac{2\sqrt2}{\sqrt K}
 \sum_{j\ge0}(2/3)^{j/2}
 =\frac{C_{\mathrm{sep}}}{\sqrt K}.
\]
This also covers vertices that are themselves deleted at a separator.
Summing over the original vertices proves the bound.  Components
already of size at most $K$ incur no charge.
\end{proof}

\subsection{Extracting a forbidden bounded pattern}

\begin{proposition}[Absence of reduced spherical arrangements]
\label{prop:no-arrangements}
With probability $1-o(1)$, $\Gamma$ admits no reduced spherical
arrangement.  In particular, for all sufficiently large admissible $n$
there exist graphs of the prescribed types, with girth at least $L$
and the diameter bound of Lemma~\ref{lem:diameter}, admitting no such
arrangement.
\end{proposition}

\begin{proof}
Let $\varepsilon>0$ be supplied by
Proposition~\ref{prop:bounded-pattern}; it is independent of the fixed
pattern bounds.  Fix $d$ for the diameter event and then the integer $D$
from Lemma~\ref{lem:short-closure}.
Splitting at length about $UL$ will make each closure cost of at most
$DL$ a small proportion of the original length.  We choose $U$ to
control that cost, $\eta$ to control the original positions discarded
by planar separation, and $C_2$ as a cutoff for the number of comparison
intervals relative to path length.  Choose, in the indicated order,
\begin{align}
 U&\ge\max\{3,\lceil96D/\varepsilon\rceil\},
       &&\text{$U$ an integer}, \notag\\
 \eta&=\min\left\{\frac1{32U},\frac{\varepsilon}{64U}\right\},
 &K&=\max\left\{1,\left\lceil(C_{\mathrm{sep}}/\eta)^2\right\rceil\right\},
 \label{eq:planar-parameters}\\
 a&=6,\qquad C_2=32a,
 &C&=K(U+D),\qquad I=\lceil C_2C\rceil. \notag
\end{align}
Here $C_{\mathrm{sep}}$ is the constant in Lemma~\ref{lem:planar-separation}.
These choices imply
\begin{equation}
 \frac{3D}{U}\le\frac{\varepsilon}{32},\qquad
 2\eta U\le\frac1{16},\qquad
 2\eta U\le\frac{\varepsilon}{32}.
 \label{eq:planar-loss-choices}
\end{equation}
For reference, the dependency order is
\begin{equation}\label{eq:parameter-order}
 \text{types},c,\delta\ ;\quad
 \varepsilon,d\ ;\quad D\ ;\quad U\ ;\quad\eta\ ;\quad K\ ;\quad C,I\ ;\quad n.
\end{equation}
All constants preceding $n$ are fixed before $n$ tends to infinity
through multiples of $2v^2$, with $L=\lfloor c\log n\rfloor$.
In particular, the bounded-pattern estimate is applied to one fixed
triple $(K,C,I)$, not to bounds growing with $n$.

Suppose, on the diameter event, that a reduced spherical arrangement
exists.  Write $N$ for its number of boundaries and $H_0$ for their
total length.  There is an ordinary cyclically immersed boundary, so
girth gives $H_0\ge L$.

\paragraph{Splitting and closing long paths.}
Keep every path of length at most $UL$ intact.  A longer path of length
$h$ is split into
$m=\lceil h/(UL)\rceil$ consecutive integer-length segments whose
lengths differ by at most one.  For large $L$, these lengths lie
between $UL/3$ and $UL$: indeed $h/m>UL/2$, and
$\lfloor h/m\rfloor\ge UL/3$ once $UL\ge6$, while
$\lceil h/m\rceil\le UL$.  For an exceptional path use its linear
order from the prescribed initial vertex, so no segment spans the
break.  Close each segment by Lemma~\ref{lem:short-closure}, and leave
every added occurrence unpaired.

Denote the number of resulting paths by $N'$.  Each has length at most
$(U+D)L$ and contains at most $UL$ original occurrences.  Since
$m\le3h/(UL)$, the total added length $H_{\mathrm{add}}$ satisfies
\begin{equation}
 H_{\mathrm{add}}\le\frac{3D}{U}H_0.
 \label{eq:closure-mass}
\end{equation}
Every resulting path has at least $L$ original occurrences, except
possibly the single exceptional path when it was kept intact: ordinary
intact paths have that length by girth, and split segments have length
at least $UL/3\ge L$.  It follows that
\begin{equation}
 N\le N',\qquad N'L\le H_0+L\le2H_0.
 \label{eq:piece-count}
\end{equation}
If the exceptional path was long, all its pieces have been closed and
there is no remaining interval path.  Otherwise the entire exceptional
path remains, with its original prescribed initial vertex.

Apply Lemma~\ref{lem:planar-intervals} to the original positions in
these pieces.  It represents their pairing by at most $aN'$ interval
pairs, each contained in the designated pieces.  On a newly closed
path, index its original segment first and its closure afterwards;
this keeps those intervals within the linear indexing as well.

\paragraph{Deleting a small set of pieces.}
Make a multigraph whose vertices are the $N'$ pieces and whose edges
are the original pairing arcs.  This graph is planar.  To see this
locally, replace each old disk vertex by a separate vertex for each
consecutive block of its incident arc ends.  Place those vertices
inside disjoint sectors of the old disk and join each to its own block
of ends within its sector.  These connections have no crossings,
including when an original arc had both ends at the old vertex.
Closures have no incident pairing arcs and need no additional edges.

By Lemma~\ref{lem:planar-separation} and the choice of $K$, delete at
most $\eta N'$ vertices so that each remaining component, called a
\emph{cluster}, has at most $K$ pieces.  The number $M$ of deleted
original positions is at most
\begin{equation}
 M\le\eta N'UL\le2\eta U H_0.
 \label{eq:deleted-mass}
\end{equation}
Retain just those occurrence pairs whose two pieces survive.  Each
surviving original position has at most one deleted partner, so the
number of newly unpaired surviving original positions is at most $M$.
Each comparison interval lies in one piece on either side; thus an
interval pair is retained or discarded in its entirety.  No retained
pair joins different clusters.

For a cluster $Z$, let $H_Z$ be its total path length including closures,
$b_Z$ the number of its unpaired positions, and $I_Z$ its number of
retained interval pairs.  Equations~\ref{eq:closure-mass},
\ref{eq:piece-count}, and~\ref{eq:deleted-mass} give
\begin{align}
 \sum_Z H_Z&\ge(1-2\eta U)H_0, \notag\\
 \sum_Z b_Z&\le(2\eta U+3D/U)H_0, \label{eq:cluster-totals}\\
 \sum_Z I_Z&\le aN'\le 2aH_0/L. \notag
\end{align}
Also every cluster satisfies $H_Z\le K(U+D)L=CL$.

\paragraph{A positive amount of admissible length survives.}
The retained total length is at least $15H_0/16$ by
Equation~\ref{eq:planar-loss-choices}.  Clusters with
$b_Z>\varepsilon H_Z$ have total length at most
\[
 \frac1\varepsilon\sum_Z b_Z\le\frac{H_0}{16}.
\]
Clusters with $I_Z>C_2H_Z/L$ likewise have total length at most
\[
 \frac{L}{C_2}\sum_Z I_Z
 \le\frac{2a}{C_2}H_0=\frac{H_0}{16}.
\]
After excluding both classes, at least $13H_0/16$ of cluster length
remains.

A cluster of length less than $L$ can contain no positive-length
cyclically immersed path.  Thus it can consist only of the one intact
short exceptional path, and there is at most one such cluster.  If its
length is $h<L$, the original arrangement also contained an ordinary
cyclic path of length at least $L$, so $H_0\ge h+L>2h$.  Removing that
cluster therefore costs less than $H_0/2$.  The clusters left after
all three exclusions have total length at least
\[
 \left(\frac{15}{16}-\frac1{16}-\frac1{16}-\frac12\right)H_0
 =\frac5{16}H_0>0.
\]
Choose one such cluster.  It contains at most $K$ positive-length
paths, has
\[
 L\le H_Z\le CL,\qquad
 b_Z\le\varepsilon H_Z,\qquad
 I_Z\le C_2H_Z/L\le C_2C\le I,
\]
and inherits Equation~\ref{eq:edge-reduction}.  All paths are cyclically
immersed except possibly the entire short exceptional interval, whose
initial vertex is still $x_A$ or $x_B$.  Added closures may traverse old
edges repeatedly; all their occurrences are unpaired, as permitted
by Proposition~\ref{prop:bounded-pattern}.  The cluster is therefore
one of the bounded patterns excluded by \Cref{prop:bounded-pattern}.

The constants in Equation~\ref{eq:planar-parameters} were chosen using
the single $\varepsilon$ independent of $K,C,I$.  The probability of
such a bounded pattern is $o(1)$, and the probability of failure of the
diameter event is $o(1)$.  Every finite reduced spherical arrangement,
regardless of size, deterministically yields a pattern for these same
fixed bounds.  This proves the claimed probability estimate and hence
the existence assertion.
\end{proof}
